\section*{Hoofdstuk \ref{ch:b2:cell-differentiation} --- Celdifferentiatie: de skeletspiercel}

\begin{solution}{exo:b2:cell-differentiation:1}
\omterm{def:b2:cell-differentiation:differentiation}{Differentiatie}: de stabiele expressie van de genen en functies van één
celtype. \omterm{def:b2:cell-differentiation:differentiation}{Determinatie}: vastlegging op een lot voordat het tot uiting komt.
\omterm{def:b2:cell-differentiation:differentiation}{Potentie}: het bereik aan lotbestemmingen dat nog openstaat. Zygote:
totipotent; cel van de \omterm{def:b2:mammal-reproduction:implantation}{binnenste celmassa}: pluripotent; satellietcel:
unipotent (alleen spier) --- toch een \omterm{def:b2:cell-differentiation:differentiation}{stamcel}; \omterm{def:b2:cell-differentiation:myogenesis}{spiervezel}:
gedifferentieerd, postmitotisch, zonder \omterm{def:b2:cell-differentiation:differentiation}{potentie}.
\end{solution}

\begin{solution}{exo:b2:cell-differentiation:2}
Cel van het dermomyotoom (Pax3) $\to$ geïnduceerde \omterm{def:b2:cell-differentiation:myogenesis}{myoblast} (\omterm{def:b2:cell-differentiation:myogenesis}{MyoD}, Myf5),
delend zolang er FGF is $\to$ uittreden uit de cyclus, myogenine $\to$
op een rij gaan staan en versmelten tot een \omterm{def:b2:cell-differentiation:myogenesis}{myotube} $\to$ spiergenen aan,
myofibrillen gebouwd (MRF4) $\to$ rijpe vezel met kernen aan de rand, plus
rustende satellietcellen (Pax7).
\end{solution}

\begin{solution}{exo:b2:cell-differentiation:3}
Dat de kern van een gedifferentieerde cel elk gen behoudt dat nodig is om
een heel dier te maken, en dat het cytoplasma van een ei haar programma kan
terugzetten: \omterm{def:b2:cell-differentiation:differentiation}{differentiatie} is omkeerbaar en geen verlies van DNA. Het
toonde niet aan dat elke gedifferentieerde kern kan worden teruggezet (de
meeste transplantaties mislukten), noch hoe het terugzetten werkt, noch dat
de gedifferentieerde cel zelf ter plaatse van lot kon veranderen.
\end{solution}

\begin{solution}{exo:b2:cell-differentiation:4}
Z-schijven die het \omterm{def:b2:cell-differentiation:fibre}{sarcomeer} begrenzen; dunne actinefilamenten die eraan
zijn verankerd; dikke myosinefilamenten, gecentreerd op de M-lijn; de A-band
(dikke filamenten) en de I-band (alleen dunne). Bij verkorting glijden de
dunne filamenten naar de M-lijn en naderen de Z-schijven elkaar: de I-band
en de H-zone worden smaller, de A-band --- de lengte van de dikke
filamenten --- verandert niet, en geen van beide filamenten verkort.
\end{solution}

\begin{solution}{exo:b2:cell-differentiation:5}
$0.2/2.5\times 10^{-6} = \num{80000}$ sarcomeren; $0.2/25\times 10^{-6}
= \num{8000}$ kernen; één \omterm{def:b2:cell-differentiation:myogenesis}{myoblast} per kern, dus ongeveer \num{8000}
versmolten \omterm{def:b2:cell-differentiation:myogenesis}{myoblasten} (satellietcellen voegen er later meer toe).
\end{solution}

\begin{solution}{exo:b2:cell-differentiation:6}
FGF induceert Id, dat het partnereiwit bindt dat \omterm{def:b2:cell-differentiation:myogenesis}{MyoD} nodig heeft om een
DNA-bindend dimeer te vormen; \omterm{def:b2:cell-differentiation:myogenesis}{MyoD} is aanwezig maar werkloos, en de cellen
blijven delen. Wanneer FGF wegvalt, wordt Id afgebroken, dimeriseert \omterm{def:b2:cell-differentiation:myogenesis}{MyoD},
wordt myogenine aangeschakeld, en treden de cellen uit de cyclus en
versmelten ze. \omterm{def:b2:cell-differentiation:myogenesis}{Myoblasten} die constitutief Id maken, versmelten nooit: ze
delen onbeperkt en vormen geen spier.
\end{solution}

\begin{solution}{exo:b2:cell-differentiation:7}
Voorwaarde $0.4m = m^{2}/(1 + m^{2}) + 0.02$. Bij $m = 2.1$: links
$0.84$, rechts $4.41/5.41 + 0.02 = 0.835$: een stationaire toestand. Lage
toestand: voor kleine $m$ is $0.4m \approx m^{2} + 0.02$, dus $m \approx
0.055$ (links $0.022$, rechts $0.023$). Verhoogd tot 1, wat boven de
drempel bij ongeveer $0.41$ ligt, klimt de cel naar $2.1$ en blijft daar:
ze is vastgelegd.
\end{solution}

\begin{solution}{exo:b2:cell-differentiation:8}
Met $k = 0.6$: bij $m = 1$ overtreft de afbraak $0.6$ de synthese
$0.52$; bij $m = 0.5$ is $0.30 > 0.22$; bij $m = 2$ is $1.2 > 0.82$: de
rechte ligt voorbij de lage toestand overal boven de kromme, en die lage
toestand is de enige stationaire toestand. Een cel waarvan \omterm{def:b2:cell-differentiation:myogenesis}{MyoD} te snel
wordt afgebroken, kan de aan-toestand niet vasthouden: ze valt terug in de
ongedifferentieerde toestand, wat de inductie ook is --- geen spier.
\end{solution}

\begin{solution}{exo:b2:cell-differentiation:9}
(a) Geen \omterm{def:b2:cell-differentiation:myogenesis}{myoblasten}, geen skeletspier: de twee factoren zijn voor de
\omterm{def:b2:cell-differentiation:differentiation}{determinatie} redundant, dus moeten ze beide weg. (b) \omterm{def:b2:cell-differentiation:myogenesis}{Myoblasten} ontstaan
en gaan op een rij staan, maar versmelten nooit en maken geen spiereiwitten:
myogenine is nodig voor de \omterm{def:b2:cell-differentiation:differentiation}{differentiatie}. (c) Bijna normaal, omdat Myf5
\omterm{def:b2:cell-differentiation:myogenesis}{MyoD} bij de \omterm{def:b2:cell-differentiation:differentiation}{determinatie} vervangt.
\end{solution}

\begin{solution}{exo:b2:cell-differentiation:10}
In een fibroblast liggen de spiergenen in chromatine dat \omterm{def:b2:cell-differentiation:myogenesis}{MyoD} kan openen;
in een levercel zijn ze misschien in heterochromatine verpakt, of
onderdrukken de eigen \omterm{prop:b2:cell-differentiation:transcriptional}{hoofdregulatoren} van de lever ze, zodat de cel niet
competent is. Toets: behandel levercellen vóór het toevoegen van \omterm{def:b2:cell-differentiation:myogenesis}{MyoD} met
een middel dat het chromatine losmaakt (een histondeacetylaseremmer), of
voeg \omterm{def:b2:cell-differentiation:myogenesis}{MyoD} toe samen met een factor die het leverprogramma verdringt, en
meet de expressie van myosine.
\end{solution}

\begin{solution}{exo:b2:cell-differentiation:11}
Aanhoudend vuren met lage frequentie verhoogt het intracellulaire calcium
langdurig; door calcium geactiveerde routes schakelen de genen van
langzaam myosine, mitochondriale biogenese, myoglobine en haarvatgroei aan,
en de snelle isovormen uit. Het hoofdprogramma van de vezel (de
MyoD-toestand, de architectuur van de sarcomeren) blijft onaangeroerd:
training verandert welke spiergenen tot expressie komen, niet de identiteit
van de cel, en geen signaal uit een zenuw kan het kraakbeenprogramma
heropenen dat bij de \omterm{def:b2:cell-differentiation:differentiation}{determinatie} werd gesloten.
\end{solution}

\begin{solution}{exo:b2:cell-differentiation:12}
De aan-toestand van een zichzelf activerende factor is een stabiel punt van
een terugkoppelingslus, in stand gehouden door voortdurende synthese en niet
door een verandering in de genen; chromatinemarkeringen voegen een tweede
laag toe die uitgeschakelde genen gesloten houdt. Terugzetten vereist dus
dat de schakeling onder haar drempel wordt gebracht en tegelijk het
chromatine wordt heropend, wat het cytoplasma van een ei of vier factoren
maar in een kleine fractie van de cellen klaarspelen: mogelijk, omdat er
niets verloren is, maar inefficiënt, omdat twee onafhankelijke sloten
tegelijk moeten worden opengebroken.
\end{solution}

\begin{solution}{pb:b2:cell-differentiation:1}
\textbf{1.} $0.2/2.5\times 10^{-6} = \num{80000}$ sarcomeren;
$0.2/25\times 10^{-6} = \num{8000}$ kernen.
\textbf{2.} Ongeveer \num{8000} per vezel; $4\times 10^{9}$ voor de
spier.
\textbf{3.} $\pi(30\times 10^{-6})^{2}\times 0.2 = \qty{5.7e-10}{m^3}$
per vezel; $2.8\times 10^{-4}$ \unit{m^3}, \qty{0.28}{L}, voor de
spier.
\textbf{4.} $5.7\times 10^{-10}/8000 = \qty{7e-14}{m^3} =
\qty{70000}{\micro m^3}$: het equivalent van 35 gewone cellen per kern.
\textbf{5.} Het inwendige is volgepakt met myofibrillen die ononderbroken
van eind tot eind moeten lopen; kernen aan de rand laten het rooster
doorlopen en liggen dicht bij het membraan en zijn signalen.
\textbf{6.} $4\times 10^{9}/2000 = 2\times 10^{6}$: 21 verdubbelingen, 10.5
dagen.
\textbf{7.} $0.4m = m^{2}/(1 + m^{2}) + 0.02$. Bij $m = 0.055$: links
$0.022$, rechts $0.003 + 0.02 = 0.023$.
\textbf{8.} Bij $m = 2.1$: links $0.84$, rechts $0.815 + 0.02 = 0.835$.
\textbf{9.} Bij $m = 0.41$: links $0.164$, rechts $0.144 + 0.02 = 0.164$.
Net daaronder is de synthese kleiner dan de afbraak, zodat $m$ naar $0.055$
daalt; net daarboven overtreft de synthese de afbraak en stijgt $m$ naar
$2.1$: elke verstoring groeit --- instabiel.
\textbf{10.} $0.6 > 0.41$: de eerste cel gaat naar de aan-toestand en
differentieert; $0.3 < 0.41$: de tweede ontspant naar $0.055$ en
differentieert niet.
\textbf{11.} Met $b = 0.4$ begint de synthesekromme bij $0.4$ en blijft ze
boven de rechte $0.4m$ tot $m \approx 3.3$: het lage snijpunt is verdwenen.
De cel differentieert spontaan, zonder inductie.
\textbf{12.} Deling halveert $m$; de dochters erven $2.1/2 =
1.05$, boven de drempel $0.41$, en klimmen terug naar $2.1$: de toestand
wordt in elke dochter hersteld. De vastlegging overleeft zolang het
gehalveerde niveau boven de drempel blijft.
\textbf{13.} $4$ per mm $\times$ \qty{200}{mm} $= 800$ satellietcellen
per vezel; \qty{10}{\%} van \num{8000} kernen, \num{800}, moet worden
vervangen.
\textbf{14.} Het beschadigde segment bevat $80$ satellietcellen; $800/80 =
10$, dus 4 verdubbelingen ($\times 16$, wat \num{1280} geeft: 800
versmelten, 480 blijven over om de voorraad aan te vullen) --- \qty{72}{h},
drie dagen.
\textbf{15.} Drie dagen deling tegenover twee tot drie weken waargenomen:
de rest is het opruimen van puin door macrofagen, de activatievertraging,
de versmelting, de bouw van myofibrillen, de herinnervatie en de rijping
van het nieuwe segment.
\textbf{16.} \omterm{def:b2:cell-differentiation:myogenesis}{MyoD} (en myogenine) moeten in de geactiveerde cellen de
spiergenen weer aanschakelen; deling en \omterm{def:b2:cell-differentiation:differentiation}{differentiatie} sluiten elkaar uit
--- Id moet dalen en de cyclus stoppen voordat myogenine kan werken en
versmelting kan plaatsvinden.
\textbf{17.} $0.95^{n} = 0.5$: $n = 13.5$, ongeveer veertien verwondingen.
\textbf{18.} $800\times 0.95^{20} = 800\times 0.36 \approx 290$.
\textbf{19.} Elke samentrekking scheurt vezels, elke scheur doet een beroep
op de voorraad, en de voorraad krimpt elke keer met enkele procenten; na
jaren is hij uitgeput, worden de gescheurde vezels niet meer herbouwd, en
vullen fibroblasten de gaten met collageen.
\textbf{20.} $(80/60)^{2} = 1.78$; kernen $8000\times 1.78 = \num{14200}$:
ongeveer \num{6200} extra uit satellietcellen.
\textbf{21.} $0.28\times 1.78 = \qty{0.5}{L}$.
\textbf{22.} Veranderd: de genen voor de myosine-isovormen, de dichtheid van
mitochondriën en haarvaten, myoglobine. Onveranderd: de spieridentiteit
(de MyoD-toestand), de architectuur van de sarcomeren, de kernen en hun
genoom.
\textbf{23.} Vezels zijn postmitotisch en kunnen zich niet vermenigvuldigen;
de enige wegen zijn dikkere vezels en meer kernen per vezel uit
satellietcellen.
\textbf{24.} Training: beperkte verdikking, omdat de vezel geen kernen kan
toevoegen; verwonding: geen herstel --- de vernietigde segmenten worden door
bindweefsel vervangen, zoals bij dystrofie.
\textbf{25.} \num{8000} kernen per vezel; drempel $m \approx 0.41$;
ongeveer drie dagen deling van satellietcellen voor een herstel.
\end{solution}
